\chapter{Aceitando Como as Coisas São}

\lettrine{Q}{uantos de vocês têm praticado} hoje tentando tornar-se algo ``Eu tenho que fazer
isto\ldots{} ou se tornar-se aquilo\ldots{} ou ver-me livre de algo\ldots{} ou
tenho que fazer algo\ldots{}'' Essa compulsividade assume o controle, mesmo na
nossa prática do Dhamma. ``Esta realidade é como é'' não é uma atitude fatalista
de não nos importarmos ou sermos indiferentes, mas é uma verdadeira abertura
para como a realidade deve ser neste momento. Por exemplo, agora precisamente, a
realidade é como é e não pode ser de outra forma neste momento. É tão óbvio,
não~é?

Neste momento, independentemente de nos estarmos a sentir eufóricos ou em baixo
ou indiferentes, felizes ou deprimidos, iluminados ou totalmente iludidos, meio
iluminados, meio iludidos, três quartos iludidos, um quarto iluminados,
esperançosos ou desesperados -- esta realidade é como é. E não pode ser de outra
forma neste momento.

Como está o nosso corpo? Reparemos simplesmente que o corpo está assim. É
pesado, está agarrado ao chão, é grosseiro, fica com fome, sente calor e frio,
fica doente, às vezes sentimo-lo muito bem, às vezes terrível. Esta realidade é
como é. Os corpos humanos são assim; de modo que essa tendência de querer que
seja diferente desaparece. Não significa que não podemos tentar tornar as coisas
melhores, mas fazemos isso com base na compreensão e na sabedoria, em vez de num
desejo ignorante.

O mundo é assim e as coisas acontecem, neva e o sol aparece, as pessoas vêm e
vão, as pessoas têm mal-entendidos, as pessoas magoam-se nos sentimentos. Ficam
preguiçosas, inspiradas, deprimidas e desiludidas, coscuvilham e desiludem-se
umas às outras, há adultério, há roubo, embriaguez, vício em drogas, e há
guerras, e sempre houve.

Aqui, numa comunidade como Amaravati, podemos ver como a realidade é. É agora
fim de semana e vêm mais pessoas oferecer comida, há mais pessoas e mais
barulho, e às vezes há crianças a gritar e a correr para cima e para baixo e
pessoas a preparar vegetais e a cortar coisas e todo um rebuliço de coisas a
acontecer ao mesmo tempo. Podemos observar ``esta realidade é como é'' em vez de
``estas pessoas estão a interferir no meu silêncio''. ``Não quero que seja
assim, quero que seja de outra forma'', pode ser a reacção se gostamos da ordem
silenciosa da refeição onde nada disso está a acontecer e onde não há barulho ou
sons agressivos. Mas a vida é assim, a vida é como é, isto é a existência
humana. Então, nas nossas mentes, abraçamos tudo isto, e ``A realidade é como
é'' permite-nos aceitar as mudanças e movimentos do silencioso ao barulhento, do
controlado e ordenado ao confuso e desordenado.

Podemos ser budistas muito egoístas e desejarmos que a vida seja muito tranquila
e sermos capazes de ``praticar'' e ter muito tempo para sentar, muito tempo para
estudar o Dhamma e pensar ``Eu não quero receber convidados e falar com as
pessoas sobre coisas estúpidas'' e ``Não quero\ldots{} blá blá blá.'' É possível
uma pessoa ser muito, mas muito mesmo egoísta, como monge budista. Pode-se
querer que o mundo se alinhe com nossos sonhos e ideais e, quando isso não
acontece, já não o querer mais. Mas, em vez de querer fazer as coisas à nossa
maneira, a via do Buddha é perceber como as coisas são. E é um grande alívio
quando aceitamos a realidade como é, mesmo que não seja muito simpática; porque
a única miséria real é não querermos que seja assim.

Quer as coisas estejam a ir bem ou não tão bem, se não estamos a aceitar as
coisas como são, então a mente tem tendência para criar uma qualquer forma de
miséria. E, se estamos apegados ao facto das coisas correrem bem, vamos começar
a preocupar-nos com elas se não correrem tão bem, mesmo quando as coisas estão
realmente a correr bem. Acabei de perceber que com pequenas coisas, como quando
está um dia de sol e a pessoa salta de alegria -- que então o próximo pensamento
será, ``Mas sabes, em Inglaterra, o sol pode desaparecer a qualquer momento''.

Assim que entendi uma percepção e estou a saltar de alegria com o sol, logo
surge o pensamento desagradável de que pode não durar. Tudo a que estamos
apegados trará o oposto. Então, quando as coisas não estão a correr muito bem, a
mente tem tendência para pensar: ``Eu quero que as coisas fiquem melhores do que
isto''. Então, o sofrimento surge onde quer que haja este apego do desejo.

O mundo sensorial é agradável e doloroso, é bonito e feio, é neutro; encontramos
nele todas as gradações, todas as possibilidades. É disto que trata a
experiência sensorial. Mas quando estão a funcionar a ignorância e a opinião
egoísta, eu só quero prazer e não quero dor. Eu quero apenas beleza e não
fealdade. ``Por favor, meu Deus, por favor, torne-me saudável, dê-me uma boa
aparência, atractividade física e deixe-me ficar jovem por muito tempo, obter
muito dinheiro, riqueza e poder, nenhuma doença, nenhum cancro, montes de coisas
bonitas ao meu redor, a envolver-me com a beleza e os prazeres dos sentidos no
seu melhor, por favor.'' E então o medo surgirá de que talvez eu obtenha o pior.
Eu poderia contrair lepra, SIDA, doença de Parkinson ou cancro. E poderia ser
rejeitado, desprezado e humilhado, deixado sozinho ao frio, faminto, doente e em
perigo, com os lobos a uivar e o vento a soprar.

Do ponto de vista do ego, existe um medo tremendo de rejeição, ostracismo ou de
se ser desprezado na nossa sociedade. Existe o medo de ser deixado sozinho e
indesejado, existe o medo de ficar velho e morrer sozinho, existe o medo natural
do perigo físico, de estarmos em situações onde os nossos corpos estão em
perigo; e há o medo do desconhecido, do misterioso, dos fantasmas e dos
espíritos invisíveis.

Então, gravitamos para a segurança, certo? Lugares confortáveis com
electricidade, aquecimento central, com seguros e garantias em tudo -- taxas
pagas e contractos legais. Tudo isto dá-nos uma sensação de segurança ou
procuramos segurança emocional. ``Diz que sempre me amarás, meu amor. Diz que
vais amar-me, mesmo que não sintas isso. Torna tudo seguro e protegido.'' E
nessa procura irá haver sempre ansiedade por causa do apego ao desejo.

Portanto, estamos a desenvolver uma luz em torno da elevação do espírito humano,
mais do que de garantias materiais. Como mendicantes, corremos o risco de não
conseguirmos nada para comer. Podemos não ter um abrigo, podemos não ter remédio
algum realmente bom, podemos não ter nada bonito para vestir. As pessoas são
muito generosas, mas, como mendicantes, não consideramos qualquer oferta como
garantida, assumindo que a merecemos. Agradecemos tudo o que nos é oferecido e
cultivamos a atitude de querer e de precisar de pouco. Precisamos de nos
preparar para partir e abrir mão de tudo a qualquer momento, para se ter o tipo
de mente que não pensa: ``Esta é a minha casa, quero que esteja garantida para o
resto da minha vida.''

Não importa que caminho tome a vida, adaptamo-nos, as nossas necessidades são
poucas. Portanto adaptamo-nos à vida, ao tempo e ao lugar, em vez de fazermos
exigências. Seja como for, é a realidade como é.

Quaisquer que sejam as doenças que eu apanhe, ou tragédias ou catástrofes ou
sucessos, ou do melhor para o pior, podemos dizer que a realidade é como é. E
nisso existe aceitação e ausência de raiva, ausência de ganância e a capacidade
de lidar com a vida enquanto está a acontecer.

Nós não estamos aqui para nos tornarmos algo ou para nos livrarmos de algo, para
mudarmos algo ou fazer algo para nós mesmos, ou para exigir algo, mas para
despertar mais e mais, para reflectir, observar e conhecer o Dhamma. Não se
preocupem que possa mudar para pior. Seja qual for a via da mudança, temos a
sabedoria de nos adaptarmos a isso. E isso eu consigo ver, é o verdadeiro
destemor de uma vida de mendicante. Podemos adaptar-nos, podemos sabiamente
aprender de todas as condições, porque este período de vida não é o nosso
verdadeiro lar.

Este tempo de vida é uma transição em que estamos envolvidos, é uma jornada
através do reino sensorial e não há ninhos, nem casas, nem habitação neste reino
sensorial. É tudo muito impermanente, sujeito a interrupções e mudanças a
qualquer momento. Essa é a sua natureza. Essa é a realidade que é. Não há nada
de deprimente nisto se não mais exigirmos segurança nisso.

A realidade da existência é que não há nenhum lar aqui. Portanto, a vida
sem-lar, a ida para uma vida de mendicância, é aquilo a que eles chamam de
mensageiro celestial, porque o espírito religioso já não partilha as ilusões da
mente mundana, que está muito determinada a ter um lar material e segurança.
Temos a confiança no Buddha, Dhamma, Sangha e no ensinamento, e nas
oportunidades como mendicantes e meditadores para a compreensão lúcida e o
entendimento que liberta a mente das ansiedades que têm origem no apego ao reino
sensorial como lar.

A ideia de possuirmos e apegar-nos às coisas é a ilusão da vida mundana. A
perspectiva do ego envia todos esses delírios nos quais temos que nos proteger o
tempo todo. Estamos sempre em perigo, há sempre algo com que nos preocuparmos,
algo a temer. Mas quando essa ilusão é perfurada com sabedoria, então há
destemor; vemos que é uma jornada, uma transição do reino sensorial e estamos
dispostos a aprender as lições que ela nos ensina, não importando que lições
possam ser.
